% Vietnamese translation for OI Public Library
% Source-File: IOI中国国家候选队论文/2007/day2/1.杨弋《Hash在信息学竞赛中的一类应用》.ppt
\documentclass[11pt]{article}
\usepackage{oipl-vn}

\newcommand{\Oh}{\mathcal{O}}
\newcommand{\Hash}{\operatorname{Hash}}
\newcommand{\xor}{\mathbin{\mathrm{xor}}}

\title{Bản trình chiếu: Một lớp ứng dụng của hàm băm trong thi tin học}
\author{Yang Yi}
\date{2007}

\begin{document}
\maketitle

\origtitle{Hash trong một lớp ứng dụng của thi tin học}
\sourcepath{IOI China candidate team papers / 2007 / day 2 / Yang Yi.ppt}

\section{Dẫn nhập}

Slide mở đầu phân biệt hai cách dùng hàm băm. Cách quen thuộc là dùng hàm
băm để đặt phần tử vào bảng băm; nếu hàm băm phân tán tốt và bảng đủ rộng,
truy vấn có thể xem như \(\Oh(1)\). Cách thứ hai, là trọng tâm của bài nói,
là dùng giá trị băm như một chữ ký ngắn của đối tượng: ta trực tiếp so sánh
hai chữ ký để phán đoán hai xâu, hai mảng hoặc hai cấu trúc có giống nhau
hay không.

Ý tưởng này cũng xuất hiện trong MD5, CRC32 hoặc SHA-1: bản thân hàm băm có
thể dùng độc lập với bảng băm. Trong thi lập trình, ta thường chấp nhận một
xác suất va chạm rất nhỏ để đổi lấy tốc độ kiểm tra tương đương rất cao.

\section{Ví dụ 1: Ghép mẫu nhiều chiều}

Bài toán một chiều là tìm vị trí xuất hiện đầu tiên của xâu mẫu trong xâu văn
bản; KMP giải được rất gọn. Khi mở rộng lên \(k\) chiều, văn bản \(X\) có
kích thước \(N_1,N_2,\ldots,N_k\), mẫu \(Y\) có kích thước
\(M_1,M_2,\ldots,M_k\), với
\[
  N=N_1N_2\cdots N_k,\qquad M=M_1M_2\cdots M_k.
\]
Điều kiện trong slide là \(k\le 10\), \(N_i\ge M_i\), và cả \(N,M\) không
vượt quá \(500000\).

KMP nhiều chiều đạt \(\Oh(k(N+M))\), nhưng khó nhớ và khó cài đúng trong thời
gian thi. Vét cạn thì liệt kê mọi điểm bắt đầu rồi so sánh toàn bộ mảng con,
tốn \(\Oh(NM)\) trong trường hợp xấu. Cách dùng băm giữ khung vét cạn ở mức
điểm bắt đầu, nhưng thay phép so sánh mảng con bằng phép so sánh giá trị băm.
Chỉ khi giá trị băm bằng nhau ta mới cần kiểm tra kỹ hơn.

Với xâu một chiều, dùng băm Rabin--Karp:
\[
  f(S)=\sum_{i=1}^{|S|} S[i]p^{|S|-i}\pmod q.
\]
Nếu \(S_i=X[i..i+M-1]\), thì khi cửa sổ trượt sang phải một ký tự:
\[
  f(S_{i+1})=\bigl[p f(S_i)+X[i+M]-p^M X[i]\bigr]\pmod q.
\]
Nếu bỏ phép lấy dư, công thức này chỉ là phép dịch trái trong hệ cơ số \(p\),
bỏ chữ số đầu và thêm chữ số cuối. Vì vậy mọi băm của các xâu con độ dài
\(M\) được tính trong thời gian tuyến tính.

\section{Băm hai chiều và nhiều chiều}

Một ma trận con \(M_1\times M_2\) có thể xem như một xâu gồm \(M_2\) dải dọc.
Trước hết tính băm cho mọi dải dọc độ dài \(M_1\), sau đó coi các giá trị ấy
là ký tự để tính băm theo chiều ngang. Hai chiều phải dùng hai cơ số khác
nhau; nếu dùng cùng một cơ số, những cấu hình như
\[
\begin{bmatrix}
a & b\\
a & a
\end{bmatrix}
\qquad\text{và}\qquad
\begin{bmatrix}
a & a\\
b & a
\end{bmatrix}
\]
có thể cho cùng kết quả băm dù ma trận không giống nhau.

Một dạng viết cho băm hai chiều là
\[
  f_2(S)=
  \sum_{i_1=1}^{\operatorname{len}_1(S)}
  \sum_{i_2=1}^{\operatorname{len}_2(S)}
  p_1^{\operatorname{len}_1(S)-i_1}
  p_2^{\operatorname{len}_2(S)-i_2}
  S[i_1,i_2]\pmod q.
\]
Khi trượt theo chiều ngang, phần thêm vào và phần bỏ ra chính là hai dải dọc
đã có băm một chiều; vì vậy một bước cập nhật vẫn là hằng số nếu các lũy thừa
đã được chuẩn bị trước.

Lên \(k\) chiều, ta lặp cùng tư tưởng: tính băm cho các mảng con
\(M_1\times\cdots\times M_{k-1}\), rồi dùng chúng như ký tự khi trượt theo
chiều thứ \(k\). Cần \(k\) lượt tính, nên độ phức tạp là \(\Oh(kN+M)\).

\section{Xác suất va chạm và chọn tham số}

Nếu hàm băm lý tưởng sinh được \(S\) giá trị khác nhau, xác suất hai đối tượng
khác nhau có cùng băm vào khoảng \(1/S\). Hàm băm trong bài không bảo đảm lý
tưởng, nhưng quy tắc thực dụng vẫn đúng: tăng không gian giá trị và dùng tham
số độc lập làm giảm rủi ro va chạm. Nếu một băm chưa đủ tin cậy, có thể tính
thêm một băm khác; đổi lại hằng số thời gian tăng rõ rệt.

Slide khuyến nghị không chọn \(p\) quá nhỏ. Với \(q\), thường chọn số nguyên
tố lớn; \(p\) nên tránh các chu kỳ ngắn modulo \(q\), ví dụ tránh để
\(p^i\equiv 1\pmod q\) với nhiều \(i\) nhỏ. Ngoài cửa sổ trượt, cùng kiểu băm
có thể được duy trì bằng chia khối, phân đoạn hoặc cây đoạn.

\section{Ví dụ 2: Equal squares}

Ural 1486 yêu cầu tìm hai hình vuông con giống nhau trong ma trận ký tự
\(N\times M\), sao cho cạnh hình vuông là lớn nhất. Cách tự nhiên là nhị phân
độ dài cạnh, rồi dùng băm hai chiều để kiểm tra xem có hai hình vuông con
cùng nội dung hay không. Nếu lưu mọi băm trong bảng băm, một lần kiểm tra
tốn \(\Oh(NM)\), tổng cộng \(\Oh(NM\log(NM))\).

Điểm đáng chú ý của slide là bài này có giới hạn thời gian chặt. Tác giả đã
chấp nhận coi hai hình vuông có cùng giá trị băm là giống nhau, không so sánh
lại toàn bộ nội dung. Để giảm nhầm lẫn do đụng vị trí trong bảng băm, có thể
tính hai giá trị: một giá trị quyết định ô trong bảng, một giá trị dùng để so
sánh với phần tử đã lưu ở ô đó.

\section{Ví dụ 3: So khớp không ổn định}

Bài toán có hai xâu \(A\) và \(B\). Xâu \(A\) bị cập nhật bởi các thao tác
chèn, xóa và đảo đoạn; truy vấn hỏi LCP của hậu tố \(A[i..]\) và \(B[j..]\).
Nếu hai xâu tĩnh, hậu tố mảng đủ để trả lời LCP, nhưng ở đây mỗi cập nhật có
thể làm thay đổi mạnh toàn bộ cấu trúc hậu tố.

Với một độ dài \(k\), ta kiểm tra hai tiền tố độ dài \(k\) bằng cách so sánh
băm. Sau đó dùng nhị phân trên \(k\) để tìm LCP. Vấn đề còn lại là duy trì
băm của đoạn con trong xâu động.

Nếu biết băm của \(S_1\) và \(S_2\), băm của phép nối là
\[
  f(S_1+S_2)=\bigl[p^{|S_2|}f(S_1)+f(S_2)\bigr]\pmod q.
\]
Vì vậy có thể dùng danh sách chia khối hoặc splay để duy trì xâu. Với chia
khối, mỗi khối lưu cả băm xuôi và băm ngược để đảo một khối trong \(\Oh(1)\);
các thao tác chèn, xóa, đảo tốn khoảng \(\Oh(\sqrt n)\), còn truy vấn LCP tốn
\(\Oh(\sqrt n\log n)\). Với splay có lazy reverse, cập nhật là \(\Oh(\log n)\)
khấu hao và truy vấn LCP là \(\Oh(\log^2 n)\).

\section{Ví dụ 4: Phán định đẳng cấu}

Để so sánh hai cây có gốc, ta tính băm từ dưới lên. Với đỉnh \(v\), lấy băm
của các con, sắp xếp tăng dần để bỏ phụ thuộc vào thứ tự con, rồi trộn chúng
bằng một hàm đủ mạnh. Việc sắp xếp là bắt buộc: cùng một cây có thể được lưu
với thứ tự con khác nhau.

Slide nhấn mạnh sự khác nhau giữa sai số ngẫu nhiên và sai số hệ thống. Một
va chạm ngẫu nhiên có thể giảm bằng cách đổi \(p,q\) hoặc dùng thêm băm thứ
hai; còn một công thức băm không phù hợp có thể tạo phản ví dụ có cấu trúc,
và đổi hằng số chưa chắc cứu được.

Với cây vô căn, ta có thể tính giá trị \(F(a,b)\) cho mỗi cạnh có hướng:
\(F(a,b)\) biểu diễn băm của cây con gốc \(b\) khi coi \(a\) là cha. Một DFS
tính mọi hướng đi xuống; một DFS đổi gốc tính các hướng đi lên bằng tiền tố
và hậu tố trên danh sách băm đã sắp xếp của các láng giềng. Sau khi mọi
\(F(a,b)\) đã biết, băm khi chọn từng đỉnh làm gốc được tính nhanh, rồi so
sánh hai cây bằng cách so khớp tập các băm đỉnh.

\section{Ví dụ 5 và 6}

Với bài biểu thức tương đương NOIP 2005, cách chuẩn là khai triển đa thức và
gộp hạng tử đồng dạng. Nhưng những biểu thức như lũy thừa lồng rất cao không
thể khai triển thực tế. Cách băm là gán vài giá trị cho biến \(a\), tính giá
trị biểu thức modulo một số lớn, rồi so sánh các kết quả. Đây cũng là một
hàm băm đặc biệt cho biểu thức.

Trong bài đoán số \(N^N=S\) với \(S\) rất dài và có thể sai vài chữ số, phần
quan trọng vẫn là tránh dựng số khổng lồ theo cách trực tiếp. Ta dùng độ dài,
ước lượng ứng viên \(N\), rồi kiểm chứng bằng các phép tính băm hoặc modulo
khác nhau. Tư tưởng chung của toàn bộ deck là: khi đối tượng gốc quá lớn để
so sánh trực tiếp, hãy chọn một chữ ký rẻ, đủ phân biệt và dễ cập nhật.

\end{document}
